\section{Fixed collision coefficients and the arithmetic return local law}
\label{sec:v40-fixed-count-local}
We now use long powers, rather than clock averages. There are two
statements with different parameter quantifiers. The finite spectral
reduction is uniform over the entire radius interval. The explicit
arithmetic local law is proved at each fixed radius. We do not assume
that a measurable phase, or its residue, varies continuously through
a parameter at which an arithmetic component appears or disappears.

\subsection{A uniform finite spectral reduction of the exact inverse}
For a roof test $q\in L^1(\R)$ with $\widehat q\in L^1(\R)$ supported
in $[-B,B]$, retain $P^q_{n,m,R}$ from
\eqref{eq:v38-band-return-functional}.

\begin{theorem}[Fixed-count resonant reduction, uniform in the radius]
\label{thm:v40-uniform-resonant-reduction}
For each fixed $B$ there are finitely many local parameter-frequency
charts, rank-one spectral branches $\lambda_j(R,\theta),\Pi_j(R,\theta)$,
and a subordinate partition of unity $\chi_0,\chi_1,\ldots,\chi_J$
on $I\times\T^2\times[-B,B]\times\T$, with the following properties.
The support of $\chi_0$ is a compact nonresonant set. Each other chart
is centered at a genuine measurable resonance; its isolated branch
is retained throughout that chart even at nearby parameters where
its modulus is less than one. Put
\[
 A_j(R,\theta)=\ell(M_{\eta_R}\Pi_j(R,\theta)(\eta_R\nu)).
\]
There are $C_B<\infty$ and $\rho_B<1$ such that, uniformly in
$R,n,m,k,t$,
\begin{align}
 P^q_{n,m,R}(k,t)
  ={}&\frac1{c(2\pi)^4}\sum_{j=1}^J
       \int_{\T^2}\int_{-B}^B\int_\T
       e^{-i(u\cdot k+bt+vn)}\widehat q(-b)
       \chi_j(R,\theta)\lambda_j(R,\theta)^m A_j(R,\theta)
                         \,dv\,db\,du\notag\\
     &+O(C_B\|\widehat q\|_1\rho_B^m).
                                      \label{eq:v40-fixed-spectral-inverse}
\end{align}
At an actual resonance the amplitude and quadratic expansion are
\eqref{eq:v40-endpoint-residue} and
\eqref{eq:v40-resonant-expansion}. The old remainder
\eqref{eq:v38-fixed-torus-remainder} has the same reduction with each
integrand multiplied by $1-\chi_B(v)$, with the same exponential
error. In particular no continuous near-unit spectral remainder is
left outside these finitely many branches on a fixed band.
\end{theorem}
\begin{proof}
At each nonresonant point Theorem~\ref{thm:v40-full-peripheral} gives
a local decaying contour. At each resonant point its simple
peripheral eigenvalue has a local rank-one continuation and an
exponentially decaying complementary contour. Strong-to-weak spectral
stability supplies uniform bounds in a smaller neighborhood. Cover
the compact joint parameter-frequency set by finitely many such
neighborhoods, and choose a partition subordinate to smaller sets
with compact closure in them. Combine nonresonant members as
$\chi_0$. The remaining $\chi_j$ may be taken continuous in the
parameter and smooth in the frequency coordinates. There is no need
to select a globally continuous generator of a resonance group.

The exact physical pairing of each power is
\[
 \int\eta_R(x)\eta_R(T_R^m x)
        e^{i(u\cdot K_{m,R}^{\rm c}+bS_{m,R}+vA_{m,R})}\,d\nu(x).
\]
Insert the local spectral splitting after multiplying by its partition
function. The error in this pairing is bounded by $C_B\rho_B^m$,
since the source and terminal section multiplier norms are uniformly
bounded. The exact return disintegration, the three torus inversions
and the roof Fourier inversion now give
\eqref{eq:v40-fixed-spectral-inverse}. Its error is at most the
uniform pairing error times the finite torus volume and
$\|\widehat q\|_1$. Multiplication by the fixed $1-\chi_B$ proves
the remainder version. All statements concern the power at the
specified $m$; neither Cesaro nor Abel averaging is used here.
\end{proof}

\subsection{The finite arithmetic coefficient}
Fix a radius $R$. By Theorem~\ref{thm:v40-no-roof-resonance} choose
the unique generator whose occupation coordinate is $2\pi/d$, where
$d=d_R$. Write its other coordinates as
\[
 u=2\pi a/d,
 \qquad s=2\pi h/d,
 \qquad a\in(\Z/d\Z)^2,\quad h\in\Z/d\Z.
\]
Its circle-valued transfer function can be normalized so that $q_*^d=1$:
indeed the $d$th power is invariant and hence constant. Write
\begin{equation}\label{eq:v40-section-phase-masses}
 w_l=\nu\{x\in Y_R^*:q_*(x)=e^{2\pi il/d}\},
       \quad l\in\Z/d\Z,\qquad \sum_lw_l=c.
\end{equation}
For $d=1$ take $a=h=0$, $q_*=1$ and $w_0=c$.
Define
\begin{equation}\label{eq:v40-arithmetic-factor}
 r_R(k,n,m)=a\cdot k+n+hm\pmod d,
 \qquad
 \mathfrak a_R(k,n,m)=\frac d{c^2}\sum_{l\bmod d}w_lw_{l+r_R(k,n,m)}.
\end{equation}
Changing the normalization of $q_*$ only permutes the labels and does
not change this factor. It is nonnegative, at most $d$, and its
average over any $d$ consecutive return indices is one.

\begin{theorem}[The exact-index, fixed-interval return law]
\label{thm:v40-arithmetic-return-LLT}
Fix $R\in I$, a bounded interval $J$ of positive length, and $M<\infty$.
Then, uniformly in the integer targets $k,n,m$ and real $t$ with
\[
 Z=\left(k/\sqrt m,(t-m\bar\tau_R)/\sqrt m,(n-mc)/\sqrt m\right),
                        \qquad |Z|\le M,
\]
one has
\begin{equation}\label{eq:v40-arithmetic-interval-LLT}
 m^2P_{n,m,R}(k,t;J)
       =c|J|\mathfrak a_R(k,n,m)g_{\Omega_R}(Z)+o_R(1).
\end{equation}
Equivalently, for
$V_n=n^{-1/2}(k_1,k_2,m-n/c,t-n\bar\tau_R/c)$,
\begin{equation}\label{eq:v40-arithmetic-original-normalization}
 n^2P_{n,m,R}(k,t;J)
       =|J|\mathfrak a_R(k,n,m)g_{D_R}(V_n)+o_R(1).
\end{equation}
If $\mathfrak a_R(k,n,m)=0$, this event has probability exactly zero,
for every roof interval. On its positive residue classes the central
probabilities are bounded below by $c_{R,M,J}m^{-2}$ for all
sufficiently large $m$.
The asserted $o_R(1)$ is at fixed radius; a global radius-uniform
arithmetic asymptotic and a pointwise roof density are not included.
\end{theorem}
\begin{proof}
Begin with a fixed integrable roof test $q$ with integrable compactly
supported Fourier transform. At this fixed radius there are only
$d$ resonances, all with $b=0$. Choose disjoint small frequency
neighborhoods of them. Outside these neighborhoods the full-band
powers decay exponentially by Theorem~\ref{thm:v40-full-peripheral}.
Inside the neighborhood of $\gamma$, use
\eqref{eq:v40-resonant-expansion}. The complementary powers decay
exponentially. With $\delta=y/\sqrt m$, the inverse phase and the
linear eigenvalue phase cancel their means and leave $e^{-iy\cdot Z}$.
The constant phase is
\[
 e^{-i(u_\gamma\cdot k+v_\gamma n+ms_\gamma)}.
\]
The projector amplitude tends to
$|\nu(\eta_Rq_\gamma)|^2$, and
$\widehat q(-y_3/\sqrt m)\to\int q$.
The Gaussian upper bound dominates the rescaled integrand; cubic
remainders tend to zero on each compact $y$ set. Splitting the
rescaled integral into a fixed ball and its Gaussian tail proves,
uniformly on the target compact set,
\begin{align}\label{eq:v40-band-arithmetic-LLT}
 m^2P^q_{n,m,R}(k,t)
   =\frac{\int q}{c}g_{\Omega_R}(Z)
       \sum_{\gamma\in\mathcal G_R}
       e^{-i(u_\gamma\cdot k+v_\gamma n+ms_\gamma)}
                   |\nu(\eta_Rq_\gamma)|^2+o_R(1).
\end{align}
The argument is an estimate of a fixed coefficient, not a Tauberian
passage from an average.

The transfer functions of the cyclic group can be chosen as
$q_*^j$, $0\le j<d$. Finite Fourier orthogonality gives
\begin{align}\label{eq:v40-finite-residue-sum}
 \sum_{j=0}^{d-1}e^{-2\pi ijr/d}
       \left|\sum_lw_l e^{2\pi ijl/d}\right|^2
       =d\sum_lw_lw_{l+r}=c^2\mathfrak a_R(k,n,m).
\end{align}
This also proves nonnegativity and the claimed residue average.
Equation~\eqref{eq:v40-band-arithmetic-LLT} is thus the desired
formula with $|J|$ replaced by $\int q$.

For the sharp interval use the pointwise upper and lower Fourier
envelopes already used in Theorem~\ref{thm:two-endpoint-collision-LLT}.
Their integrals are $|J|\pm2\pi/B$; the minorant need not be
nonnegative. Multiply their inequalities by
the unchanged positive exact-index source. For each fixed $B$ take
$m\to\infty$, and only then let $B\to\infty$. Since all resonances
have roof coordinate zero, the same arithmetic factor occurs on
both sides for every $B$. This proves
\eqref{eq:v40-arithmetic-interval-LLT}. In the original covariance
coordinates use \eqref{eq:v37-return-gaussian-agreement} and
$n/m=c+O(m^{-1/2})$, exactly as in
Corollary~\ref{cor:v38-return-normalization}. The bounded factor
$\mathfrak a_R\le d$ preserves the error and gives
\eqref{eq:v40-arithmetic-original-normalization}.

Finally the iterated phase equation on the exact section event gives
\[
 q_*(T_R^m x)=e^{2\pi ir_R(k,n,m)/d}q_*(x).
\]
If the sum of products
in \eqref{eq:v40-arithmetic-factor} is zero, every permitted pair
of section phase classes has a null initial or null terminal class.
Invariance of $T_R$ makes the event null. There are only finitely many
positive residue factors at fixed $R$, so their minimum is positive.
Uniform positivity of the Gaussian on the stated target compact
set gives the last lower bound.
\end{proof}

\begin{corollary}[Which exact-index complement remains]
\label{cor:v40-explicit-fixed-complement}
For every fixed radius and compactly band-limited roof test,
\[
 m^2\mathcal R^q_{n,m,R}(k,t)
   =c\left(\int q\right)
       (\mathfrak a_R(k,n,m)-1)g_{\Omega_R}(Z)+o_R(1).
\]
In particular the unmodulated fixed-interval local law follows if all
nontrivial section residues vanish, and in particular if $d_R=1$.
The general theorem does not assume either statement.
\end{corollary}
\begin{proof}
Subtract Proposition~\ref{prop:v38-exact-index-remainder} from
\eqref{eq:v40-band-arithmetic-LLT}, using
\eqref{eq:v40-finite-residue-sum}. No inference from the pure
occupation Abel limit enters this subtraction.
\end{proof}

\begin{corollary}[Selected endpoints and the fine-window limit]
\label{cor:v40-selected-arithmetic}
For nonnegative bounded endpoints $a_0,a_1$ with
$\eta_Ra_0,\eta_Ra_1$ multiplier-bounded, define
$w_l^{(i)}=\nu(\eta_Ra_i\one_{\{q_*=e^{2\pi il/d}\}})$ for $i=0,1$.
The fixed-radius selected version of
\eqref{eq:v40-arithmetic-interval-LLT} is obtained by
replacing $\mathfrak a_R$ by
\[
 \mathfrak a_R^{a_0,a_1}(k,n,m)
       =\frac d{c^2}\sum_lw_l^{(0)}w_{l+r_R(k,n,m)}^{(1)}.
\]
Its average over $d$ consecutive return indices is
$\nu_R^*(a_0)\nu_R^*(a_1)$. Thus averaging any diverging subdiffusive
block recovers the unmodulated amplitude in
Corollary~\ref{cor:v38-selected-return-window} at fixed $R$.
\end{corollary}
\begin{proof}
Use the general residue formula \eqref{eq:v40-endpoint-residue} in
the same fixed-band proof. Its finite Fourier sum equals the displayed
cross correlation of phase-class masses. The sum over all residues
factors as $(\sum_lw_l^{(0)})(\sum_lw_l^{(1)})$. Within a subdiffusive block
the Gaussian varies by $o(1)$, and the incomplete final residue block
has relative size at most $d/H_m\to0$. The error in the fixed-index
formula is uniform over the target compact set and remains $o_R(1)$
after division by the block width. No Gaussian amplitude for an
arbitrary path-dependent selector is asserted.
\end{proof}

\paragraph{The remaining raw and parameter questions.}
The full occupation-torus complement is now a controlled spectral
problem on each fixed roof band. At a fixed radius its exact-index
contribution is the finite arithmetic expression above. This does
not assert that the arithmetic factor equals one, or that a pointwise
selection of exact resonances gives a radius-uniform asymptotic near
a change of resonance type. The uniform statement in such regions is
\eqref{eq:v40-fixed-spectral-inverse}, which retains all nearby
spectral branches. Establishing triviality of the section residues,
or tracking their parameter transitions, is a distinct remaining task.

A fixed positive roof interval is not a density value. None of the
arguments here bounds the complete second-derivative coarea sum,
removes every critical edge, or permits $B$ to grow with $m$ inside
an unquantified spectral estimate. The common pointwise raw correction
and the full roof-frequency estimate in Theorem~\ref{thm:LLT} remain
requirements of that unchanged theorem. All its inherited geometry
and inversion arguments are retained.
